\chapter{Design}
\label{chap:design}

The prototype keeps the engine alive between runs and saves its state at
regular points in the document.

\section{Snapshots}
\label{sec:snapshots}

A snapshot records the engine state at a page boundary. When a source file
changes, the system finds the last snapshot taken before the first changed
byte and resumes from there. Everything before that point is reused.

\subsection{Choosing Snapshot Points}
\label{sec:snapshot-points}

Snapshots cost memory, so the system does not take one at every page. With a
snapshot every $k$ pages, the expected number of pages to retypeset after an
edit on a uniformly random page is
\begin{equation}
  E[\text{pages}] = \frac{k - 1}{2} + 1 .
  \label{eq:expected-pages}
\end{equation}

\section{Invalidation}
\label{sec:invalidation}

A change to an auxiliary file invalidates every snapshot, because any page
may read it (Section~\ref{sec:aux}). The prototype therefore compares the new
auxiliary file with the old one and restarts from the beginning only when
they differ.

\section{Limitations}
\label{sec:limitations}

The design assumes that a run is deterministic given its inputs. Documents
that read the clock or call external programs break that assumption, and the
prototype falls back to a full run for them.
